\section{Related Work and Positioning}

This paper sits in computational social choice, but its contribution is framed as a small-scope systems case study rather than as a new impossibility theorem or a general scaling result.
The object of study is the sen24 feasibility boundary: which axiom-lever combinations are SAT or UNSAT, which boundary cases admit a concrete explanation, and which artifacts can be replayed and audited afterward.
That emphasis on boundary identification, explicit evidence, and replayable outputs is what turns the case study into an AI/Systems contribution.

Constraint Hierarchies and Assumption-Based Truth Maintenance Systems are natural reference points because they both reason about conflicting constraints or assumptions.
We differ from them in output and trust model.
Our primary deliverable is not only a consistent assignment or a justification structure, but an auditable atlas whose cases are tied to concrete artifacts: frontier summaries, one MUS and one small MCS candidate for boundary explanation, committed proof-carrying UNSAT references, and validated SAT rule cards.
In other words, we use conflicting-assumption structure to organize cause and repair, but we package the result as proof-/witness-carrying evidence rather than as inference bookkeeping alone.

MAXSAT and OMT are also close neighbors because they can optimize soft-constraint violations or compute minimum relaxations.
We do not position this paper as a replacement for those methods or as a speed comparison against them.
Instead, the paper uses optimization-style reasoning only as a triangulation baseline for the stronger C6 stage: repair enumeration is compared against an independent optimum baseline, while the main contribution remains causal boundary mapping, MUS/MCS-style explanation at the paper's conservative claim level, and validated SAT witness artifacts.
The difference matters because an optimum value alone does not provide the paper's central outputs: a reproducible frontier, explicit boundary explanations, and audit-ready rule artifacts.

The systems angle is therefore not raw solving power but auditable integration.
The repository fixes a chain from semantic specification to DIMACS generation, manifest auditing, SAT/UNSAT outputs, LRAT replay in Lean for committed UNSAT references, SAT witness validation, and reproducible paper-facing artifacts.
Evidence bundles, schema/version fields, deterministic render paths, and CI-backed smoke checks are part of the contribution because they make the sen24 atlas inspectable by a reviewer without trusting the solver as the final authority.

The scope remains intentionally narrow.
We claim only a sen24 case study under a small-scope hypothesis, not a general account of arbitrary $n,m$, not a complete survey of social-choice theorem proving, and not a formal Lean proof of pruning or symmetry heuristics.
The value of the paper is that even under those limits, an auditable possibility atlas can connect boundary, cause, repair, and witness in a form that is useful for verified, AI-assisted rule design.
